\chapter{A speed-up szélsőértékei}

Ebben a fejezetben külön is megvizsgálom, hogy a teljes mérési halmazon belül
hol mértem a \textbf{legnagyobb}, illetve a \textbf{legkisebb} gyorsítást.

A speed-up definíciója a dokumentumban mindenhol ugyanaz:
\textbf{mindig a saját natív C implementációt veszem baseline-nak}, és ehhez
viszonyítom az OpenCL-es útvonalat. Vagyis a speed-up azt fejezi ki,
hogy az OpenCL-es végrehajtás hányszor gyorsabb a natív CPU-s megoldásnál.

\section{Összesített szélsőértékek}

\begin{table}[H]
  \centering
  \caption{A teljes mérési halmaz legkisebb és legnagyobb OpenCL speed-up értéke.}
  \fittable{\pgfplotstabletypeset[
    col sep=comma,
    columns={Extreme,PC,BuildArchitecture,PowerState,Algorithm,KeySizeBits,DataSizeMegabytes,Direction,mean_speedup},
    columns/Extreme/.style={string type, column name=Szélsőérték},
    columns/PC/.style={string type, column name=Gép},
    columns/BuildArchitecture/.style={string type, column name=Build},
    columns/PowerState/.style={string type, column name=Tápellátás},
    columns/Algorithm/.style={string type, column name=Algoritmus},
    columns/KeySizeBits/.style={column name=Kulcs (bit), fixed, precision=0},
    columns/DataSizeMegabytes/.style={column name=Adat (MB), fixed, precision=0},
    columns/Direction/.style={string type, column name=Irány},
    columns/mean_speedup/.style={column name=Speed-up, fixed, precision=2},
    every head row/.style={before row=\toprule, after row=\midrule},
    every last row/.style={after row=\bottomrule},
  ]{\GenPath/speedup_extremes.csv}}
\end{table}

A teljes adathalmazban a \textbf{legkisebb speed-up} a \textbf{PC3 x64, GCM-256, 64 MB, Encrypt, töltőn}
konfigurációban jelent meg, \textbf{6.95x} értékkel. A \textbf{legnagyobb speed-up} ezzel szemben a
\textbf{PC4 x86, GCM-256, 256 MB, Encrypt} esetben jelent meg, \textbf{38.97x} értékkel.

Ezt én úgy értelmezem, hogy a minimumnál egy laptopos, visszafogottabb platformot látok,
olyan esettel, ahol az OpenCL ugyan még mindig jelentősen gyorsabb a natív C baseline-nál,
de a különbség már nem extrém. A maximumnál viszont egyszerre több tényező erősíti egymást:
\begin{itemize}
  \item erős GPU (PC4),
  \item nagyobb bemenet (256 MB), tehát az OpenCL fix overhead jobban amortizálódik,
  \item x86 build, ahol a natív CPU baseline lassabb, ezért ugyanaz az OpenCL teljesítmény nagyobb speed-up-ként jelenik meg.
\end{itemize}

\section{Szélsőértékek architektúránként}

\begin{table}[H]
  \centering
  \caption{x64 és x86 szerinti legkisebb/legnagyobb speed-up értékek.}
  \fittable{\pgfplotstabletypeset[
    col sep=comma,
    columns={Extreme,Architecture,PC,PowerState,Algorithm,KeySizeBits,DataSizeMegabytes,Direction,mean_speedup},
    columns/Extreme/.style={string type, column name=Szélsőérték},
    columns/Architecture/.style={string type, column name=Arch.},
    columns/PC/.style={string type, column name=Gép},
    columns/PowerState/.style={string type, column name=Tápellátás},
    columns/Algorithm/.style={string type, column name=Algoritmus},
    columns/KeySizeBits/.style={column name=Kulcs (bit), fixed, precision=0},
    columns/DataSizeMegabytes/.style={column name=Adat (MB), fixed, precision=0},
    columns/Direction/.style={string type, column name=Irány},
    columns/mean_speedup/.style={column name=Speed-up, fixed, precision=2},
    every head row/.style={before row=\toprule, after row=\midrule},
    every last row/.style={after row=\bottomrule},
  ]{\GenPath/speedup_arch_extremes.csv}}
\end{table}

Ez a bontás még jobban alátámasztja a korábbi megfigyelést: \textbf{x86-on általában nagyobb a speed-up szám},
de ez nem azt jelenti, hogy az OpenCL kernel önmagában gyorsabb lenne x86 alatt. A fő ok inkább az,
hogy a \textbf{natív C baseline x86 buildben lassabb}, ezért ugyanaz a GPU-s teljesítmény nagyobb arányt ad.

\section{A legalacsonyabb és legmagasabb speed-up tartományok}

\subsection{Az 5 legkisebb speed-up érték}

\begin{table}[H]
  \centering
  \caption{Az 5 legkisebb OpenCL speed-up érték a teljes mérési halmazban.}
  \fittable{\pgfplotstabletypeset[
    col sep=comma,
    columns={PC,BuildArchitecture,PowerState,Algorithm,KeySizeBits,DataSizeMegabytes,Direction,mean_speedup},
    columns/PC/.style={string type, column name=Gép},
    columns/BuildArchitecture/.style={string type, column name=Build},
    columns/PowerState/.style={string type, column name=Tápellátás},
    columns/Algorithm/.style={string type, column name=Algoritmus},
    columns/KeySizeBits/.style={column name=Kulcs (bit), fixed, precision=0},
    columns/DataSizeMegabytes/.style={column name=Adat (MB), fixed, precision=0},
    columns/Direction/.style={string type, column name=Irány},
    columns/mean_speedup/.style={column name=Speed-up, fixed, precision=2},
    every head row/.style={before row=\toprule, after row=\midrule},
    every last row/.style={after row=\bottomrule},
  ]{\GenPath/speedup_bottom5.csv}}
\end{table}

Az alsó tartományban főleg \textbf{PC3} szerepel, vagyis a laptopos platform. Ez összhangban van a
korábbi tápellátásos és architektúrás vizsgálattal: itt a GPU-platform szerényebb, és a rendszer teljesítménye
érzékenyebb a futási környezetre.

\subsection{Az 5 legnagyobb speed-up érték}

\begin{table}[H]
  \centering
  \caption{Az 5 legnagyobb OpenCL speed-up érték a teljes mérési halmazban.}
  \fittable{\pgfplotstabletypeset[
    col sep=comma,
    columns={PC,BuildArchitecture,PowerState,Algorithm,KeySizeBits,DataSizeMegabytes,Direction,mean_speedup},
    columns/PC/.style={string type, column name=Gép},
    columns/BuildArchitecture/.style={string type, column name=Build},
    columns/PowerState/.style={string type, column name=Tápellátás},
    columns/Algorithm/.style={string type, column name=Algoritmus},
    columns/KeySizeBits/.style={column name=Kulcs (bit), fixed, precision=0},
    columns/DataSizeMegabytes/.style={column name=Adat (MB), fixed, precision=0},
    columns/Direction/.style={string type, column name=Irány},
    columns/mean_speedup/.style={column name=Speed-up, fixed, precision=2},
    every head row/.style={before row=\toprule, after row=\midrule},
    every last row/.style={after row=\bottomrule},
  ]{\GenPath/speedup_top5.csv}}
\end{table}

A felső tartományban szinte teljesen \textbf{PC4 x86} dominál. Ez is azt mutatja, hogy a legnagyobb
speed-up nem feltétlenül azt jelenti, hogy ott a legjobb az OpenCL abszolút teljesítménye, hanem azt,
hogy ott a \textbf{natív baseline-hoz képest} a legnagyobb a relatív előny. A dokumentum szempontjából
pont ez volt a lényeges mérőszám, mert a feladatban a saját C implementációmhoz akartam viszonyítani a
párhuzamosítással elért gyorsulást.

\cleardoublepage
